% =====================================================================
%  LiaScript / Markdown CheatSheet — A4, doppelseitig
%  Fortbildung "Interaktive OER mit LiaScript erstellen", Freiberg 2026
%  Compile:  pdflatex cheatsheet.tex   (zweimal für korrekte Boxen)
% =====================================================================
\documentclass[a4paper,landscape,9pt]{extarticle}

% --- Layout ----------------------------------------------------------
\usepackage[landscape,margin=0.9cm,top=1.0cm,bottom=0.7cm]{geometry}
\usepackage{multicol}
\setlength{\columnsep}{8pt}
\setlength{\parindent}{0pt}
\setlength{\parskip}{0pt}

% --- Sprache & Schrift ----------------------------------------------
\usepackage[T1]{fontenc}
\usepackage[utf8]{inputenc}
\usepackage[ngerman]{babel}
\usepackage{libertine}
\usepackage[varqu,varl,scaled=0.92]{inconsolata}

% --- Farben ----------------------------------------------------------
\usepackage{xcolor}
\definecolor{mdblue}{HTML}{1F4E79}
\definecolor{mdblueLt}{HTML}{E8EEF5}
\definecolor{liagreen}{HTML}{2E6B3D}
\definecolor{liagreenLt}{HTML}{E8F0E9}
\definecolor{schulred}{HTML}{A61E2D}
\definecolor{schulredLt}{HTML}{F7E9EB}
\definecolor{lootteal}{HTML}{007F84}
\definecolor{loottealLt}{HTML}{E5F4F4}
\definecolor{codebg}{HTML}{F6F6F4}
\definecolor{rulegray}{HTML}{D0D0D0}
\definecolor{muted}{HTML}{606060}

% --- Pakete ----------------------------------------------------------
\usepackage[most]{tcolorbox}
\usepackage{listings}
\usepackage{qrcode}
\usepackage{fontawesome5}
\usepackage{array}
\usepackage[normalem]{ulem}
\usepackage{hyperref}
\hypersetup{hidelinks}

% --- Listings: Markdown/LiaScript-Quellcode --------------------------
\lstdefinelanguage{lia}{
  morecomment=[l]{\#},
  morecomment=[s]{<!--}{-->},
  sensitive=true
}
\lstset{
  basicstyle=\ttfamily\scriptsize,
  breaklines=true,
  breakatwhitespace=true,
  showstringspaces=false,
  keepspaces=true,
  columns=fullflexible,
  upquote=true,
  literate=%
    {ä}{{\"a}}1 {ö}{{\"o}}1 {ü}{{\"u}}1
    {Ä}{{\"A}}1 {Ö}{{\"O}}1 {Ü}{{\"U}}1
    {ß}{{\ss}}1 {„}{{\glqq}}1 {“}{{\grqq}}1,
  frame=none,
  xleftmargin=0pt,
  aboveskip=0pt,
  belowskip=0pt
}

% --- Karten-Vorlage --------------------------------------------------
%  Zwei Varianten: mdcard (blau, Markdown) und liacard (grün, LiaScript)
%  Innen: links Quellcode, rechts Wirkung — getrennt durch \tcblower.
\tcbset{
  cheatbase/.style={
    enhanced,
    sidebyside,
    sidebyside align=top,
    lefthand width=0.46\linewidth,
    segmentation style={solid, rulegray, line width=0.3pt},
    colback=white,
    boxrule=0.3pt,
    arc=2pt,
    left=4pt,right=4pt,top=2pt,bottom=2pt,middle=4pt,
    before skip=3pt, after skip=3pt,
    breakable=false,
    fonttitle=\sffamily\bfseries\footnotesize,
    coltitle=white
  }
}

\newtcolorbox{mdcard}[1]{cheatbase, colframe=mdblue!75,  colbacktitle=mdblue,   title={#1}}
\newtcolorbox{liacard}[1]{cheatbase, colframe=liagreen!75, colbacktitle=liagreen, title={#1}}
\newtcolorbox{lootcard}[1]{cheatbase, colframe=lootteal!75, colbacktitle=lootteal, title={#1}}
\newtcolorbox{schulcard}[2]{
  cheatbase,
  colframe=schulred!75,
  colbacktitle=schulred,
  title={%
    \parbox[c]{0.47\linewidth}{#1}\hfill%
    \parbox[c]{0.50\linewidth}{\raggedleft\normalfont\ttfamily\bfseries\fontsize{6.8}{7.3}\selectfont #2}%
  }
}

% Karte für die ganze Spaltenbreite (z. B. Mini-Kurs am Ende)
\newtcolorbox{liawide}[1]{
  enhanced,
  colback=white,
  colframe=liagreen!75,
  boxrule=0.3pt, arc=2pt,
  left=4pt,right=4pt,top=2pt,bottom=2pt,
  before skip=3pt, after skip=0pt,
  breakable=false,
  fonttitle=\sffamily\bfseries\footnotesize,
  coltitle=white, colbacktitle=liagreen, title=#1
}

% Info-Box ohne sidebyside (für Tools/Templates-Spalten unten)
\tcbset{
  infobase/.style={
    enhanced, colback=white,
    boxrule=0.3pt, arc=2pt,
    left=5pt,right=5pt,top=3pt,bottom=3pt,
    before skip=3pt, after skip=0pt,
    breakable=false,
    fonttitle=\sffamily\bfseries\footnotesize, coltitle=white
  }
}
\newtcolorbox{mdinfo}[1]{infobase, colframe=mdblue!75,    colbacktitle=mdblue,   title={#1}}
\newtcolorbox{liainfo}[1]{infobase, colframe=liagreen!75, colbacktitle=liagreen, title={#1}}
\newtcolorbox{schulinfo}[1]{infobase, colframe=schulred!75, colbacktitle=schulred, title={#1}}

% Kleine Hilfsmacros für die "Wirkung"-Spalte
\newcommand{\prevtxt}[1]{{\small #1}}
\newcommand{\prevsm}[1]{{\footnotesize\color{muted} #1}}
\newcommand{\lootval}[1]{{\ttfamily\bfseries\color{lootteal}#1}}
\newcommand{\sheetfooter}{%
  \vfill
  {\sffamily\scriptsize\color{muted}
  \faInfoCircle~Live-Editor: \href{https://liascript.github.io/LiveEditor/}{\textbf{liascript.github.io/LiveEditor}}\hfill
  \faBook~Doku: \href{https://liascript.github.io}{\textbf{liascript.github.io}}\hfill
  \faGithub~\href{https://mint-the-gap.github.io/Aufgabensammlung/}{\textbf{mint-the-gap.github.io/Aufgabensammlung/}} \quad
  CC-BY-SA 4.0 — Martin Lommatzsch, Freiberg 2026}}

\pagestyle{empty}
\begin{document}

% =====================================================================
% SEITE 1 — MARKDOWN
% =====================================================================
\noindent\begin{tcolorbox}[
  enhanced, frame hidden, sharp corners,
  colback=mdblue, coltext=white,
  left=12pt, right=8pt, top=4pt, bottom=4pt,
  before skip=0pt, after skip=4pt,
  width=\linewidth,
  borderline west={4pt}{-4pt}{mdblue!55}
]
\begin{minipage}[c]{0.66\linewidth}
  {\sffamily\fontsize{26}{28}\selectfont\bfseries Markdown}\hspace{0.4em}%
  {\sffamily\large\color{white!75} — die Basis}\\[2pt]
  {\sffamily\footnotesize\itshape\color{white!90} Was jeder Markdown-Editor versteht — Grundlage für alle LiaScript-Dokumente.}
\end{minipage}\hfill
\begin{minipage}[c]{0.20\linewidth}\raggedright
  {\sffamily\scriptsize\color{white!85}\faGraduationCap~~Fortbildung~$\cdot$~Freiberg 2026}\\[2pt]
  {\sffamily\bfseries\normalsize Seite 1\,/\,5}\\[1pt]
  {\sffamily\scriptsize\color{white!75} Markdown-Grundlagen}
\end{minipage}\hfill
\begin{minipage}[c]{0.10\linewidth}\centering
  {\color{black}\colorbox{white}{\hspace{1pt}\qrcode[height=1.55cm,nolinks]{https://liascript.github.io/LiveEditor/}\hspace{1pt}}}\\[1pt]
  {\sffamily\scriptsize\color{white!90}\faEdit~Live-Editor}
\end{minipage}
\end{tcolorbox}

% --- Tag für Aufgabenbezug (klein, oben rechts in der Titelzeile) ---
\newcommand{\aufg}[1]{\hfill{\sffamily\scriptsize\color{white!85}\faTasks~A#1}}

\begingroup
\tcbset{cheatbase/.append style={top=1pt,bottom=1pt,before skip=2pt,after skip=2pt}}
\begin{multicols}{3}

% =====================================================================
% Block A — Textauszeichnung & Struktur (Aufgaben 2-4)
% =====================================================================

% --- YAML-Header -----------------------------------------------------
\begin{mdcard}{YAML-Header}
\begin{lstlisting}[language=lia,basicstyle=\ttfamily\tiny]
<!--
author:   Maria Muster
version:  1.0.0
language: de
narrator: Deutsch Female
mode:     Presentation
edit:     true
tags:     OER, LiaScript, Fach
comment:  Interaktiver Kurs
import:   ./template.md
-->
\end{lstlisting}
\tcblower
\prevsm{Steht ganz oben im Dokument.\\
Definiert Version, Sprache, Startmodus,
Bearbeitung, Tags und Beschreibung.\\[2pt]
{\footnotesize \texttt{import:} mit lokalem Pfad oder URL für jeden Baustein wiederholen.}}
\end{mdcard}

% --- R3. Version & gespeicherte Eingaben ----------------------------
\begin{mdcard}{Version \& gespeicherte Eingaben}
\begin{lstlisting}[language=lia]
<!-- version: 1.0.0 -->
<!-- version: 0.1.0 -->
<!-- version: 0.0.1 -->
\end{lstlisting}
\tcblower
\prevsm{\textbf{1.0.0:} Release; Quiz- und Codezustände werden im Browser persistent gespeichert.\\[2pt]
\textbf{0.1.0:} Entwicklungsstand; bei jedem Laden neu parsen, Eingaben nicht persistent speichern.\\[2pt]
\textbf{0.0.1:} früher Debug-/Teststand; gleiches Speicherverhalten wie jede Hauptversion~0.}
\end{mdcard}

% --- A3. Überschriften ----------------------------------------------
\begin{mdcard}{Überschriften}
\begin{lstlisting}[language=lia]
# Kapitel
## Abschnitt
### Unterabschnitt
#### Detail
\end{lstlisting}
\tcblower
{\large\bfseries Kapitel}\par\vspace{1pt}
{\normalsize\bfseries Abschnitt}\par\vspace{1pt}
{\small\bfseries Unterabschnitt}\par\vspace{1pt}
{\footnotesize\bfseries Detail}
\end{mdcard}

% --- A2a. Textauszeichnung ------------------------------------------
\begin{mdcard}{Textauszeichnung}
\begin{lstlisting}[language=lia]
**fett**, *kursiv*,
~~durchgestrichen~~,
`Quellcode`
\end{lstlisting}
\tcblower
\prevtxt{\textbf{fett}, \textit{kursiv},
\sout{durchgestrichen}, \texttt{Quellcode}}
\end{mdcard}

% --- A2b. Absätze & Trennlinie --------------------------------------
\begin{mdcard}{Absätze \& Trennlinie}
\begin{lstlisting}[language=lia]
Erster Absatz.

Zweiter Absatz nach
einer Leerzeile.

---
\end{lstlisting}
\tcblower
\prevtxt{Erster Absatz.\\[3pt]
Zweiter Absatz nach einer Leerzeile.\\[2pt]
\textcolor{rulegray}{\hrulefill}}
\end{mdcard}

% --- A4. Listen -----------------------------------------------------
\begin{mdcard}{Listen}
\begin{lstlisting}[language=lia]
- erstens
- zweitens

  - eingerückt

- drittens

1. eins
2. zwei
\end{lstlisting}
\tcblower
\prevtxt{%
$\bullet$ erstens\\
$\bullet$ zweitens\\[2pt]
\hspace*{1em}$\circ$ eingerückt\\[2pt]
$\bullet$ drittens\\[3pt]
1. eins\\
2. zwei\\[2pt]
{\scriptsize\color{muted}Untergeordnete Listen um \textbf{2 Zeichen} einrücken.}}
\end{mdcard}

% --- R1. Aufgabenliste / Checkboxen ---------------------------------
\begin{mdcard}{Aufgabenliste}
\begin{lstlisting}[language=lia]
- [x] erledigt
- [ ] offen
- [ ] noch offen
\end{lstlisting}
\tcblower
\prevtxt{%
\faCheckSquare~erledigt\\
\faSquare~offen\\
\faSquare~noch offen}
\end{mdcard}

% =====================================================================
% Block B — Tabellen, Code & Gleichungen (Aufgaben 5-6)
% =====================================================================

% --- A5. Tabellen ---------------------------------------------------
\begin{mdcard}{Tabellen}
\begin{lstlisting}[language=lia]
<!-- data-type="none"
  data-sortable="false" -->
| Operator | Wirkung      |
|:---------|:------------:|
| AND      | verkleinert  |
| OR       | vergrößert |
| NOT      | schließt aus|
\end{lstlisting}
\tcblower
\prevtxt{%
\renewcommand{\arraystretch}{1.05}%
\begin{tabular}{|l|c|}\hline
\textbf{Operator} & \textbf{Wirkung}\\\hline
AND & verkleinert\\
OR  & vergrößert\\
NOT & schließt aus\\\hline
\end{tabular}\\[2pt]
{\scriptsize\color{muted}\texttt{:--} links, \texttt{:--:} zentriert, \texttt{--:} rechts.\\
\texttt{data-type="none"}: keine Grafik; \texttt{data-sortable="false"}: feste Reihenfolge.}}
\end{mdcard}

% --- A6a. Inline-Code & Codeblock -----------------------------------
\begin{mdcard}{Inline-Code \& Codeblock}
\begin{lstlisting}[language=lia]
Operatoren wie `AND`
und `OR` als Code.

```python
print("Hallo")
```
\end{lstlisting}
\tcblower
\prevtxt{Operatoren wie \texttt{\small AND}
und \texttt{\small OR} als Code.\\[3pt]
\colorbox{codebg}{\parbox{0.85\linewidth}{\ttfamily\scriptsize print("Hallo")}}}
\end{mdcard}

% --- A6b. Mathematik ------------------------------------------------
\begin{mdcard}{Mathematische Gleichungen}
\begin{lstlisting}[language=lia]
Inline: $A \cdot B$

Block:
$$\neg A \lor (B \land C)$$
\end{lstlisting}
\tcblower
\prevtxt{Inline: $A \cdot B$\\[3pt]
Block:\\
$\displaystyle \neg A \lor (B \land C)$\\[3pt]
{\scriptsize\color{muted}\texttt{\textbackslash land}~$\land$,\,
\texttt{\textbackslash lor}~$\lor$,\,
\texttt{\textbackslash neg}~$\neg$,\,
\texttt{\textbackslash cdot}~$\cdot$}}
\end{mdcard}

% =====================================================================
% Block C — Links, Bilder, Hinweise (Aufgaben 7-8)
% =====================================================================

% --- A7a. Links ------------------------------------------------------
\begin{mdcard}{Links}
\begin{lstlisting}[language=lia]
[Beschreibung](https://...)

<https://example.org>
\end{lstlisting}
\tcblower
\prevtxt{\textcolor{mdblue}{\underline{Beschreibung}}\\[2pt]
\textcolor{mdblue}{\underline{https://example.org}}\\[3pt]
{\scriptsize\color{muted}Sprechende Beschreibung erleichtert die Orientierung.}}
\end{mdcard}

% --- R2. Fußnoten ---------------------------------------------------
\begin{mdcard}{Fußnoten}
\begin{lstlisting}[language=lia]
Wichtig[^1] zum Lernen.

[^1]: Quelle: ...
\end{lstlisting}
\tcblower
\prevtxt{Wichtig\textsuperscript{\textcolor{mdblue}{1}} zum Lernen.\\[3pt]
{\scriptsize \textsuperscript{1}~Quelle: \dots}}
\end{mdcard}

% --- A7b. Bilder -----------------------------------------------------
\begin{mdcard}{Bilder}
\begin{lstlisting}[language=lia]
![Alt-Text](pfad/bild.png)

![Venn-Diagramm AND/OR](
  https://.../venn.svg)
\end{lstlisting}
\tcblower
\prevtxt{\fbox{\rule{0pt}{1.2em}\hspace{1.2em}{\scriptsize\color{muted}Bild}\hspace{1.2em}}\\[2pt]
{\scriptsize\color{muted}\textit{Alt-Text in eckigen Klammern — wichtig für Barrierefreiheit.}}}
\end{mdcard}

% --- R5. Audio ------------------------------------------------------
\begin{mdcard}{Audio}
\begin{lstlisting}[language=lia]
?[Audio](datei.mp3)

?[Hörbeispiel](https://...
  /aufnahme.ogg)
\end{lstlisting}
\tcblower
\prevsm{\faVolumeUp~MP3, OGG und SoundCloud als Audioplayer einbetten.\\[2pt]
\textbf{?} vor dem Link kennzeichnet eine Audioressource.}
\end{mdcard}

% --- R4. Medien abspielbar ------------------------------------------
\begin{mdcard}{Medien abspielbar}
\begin{lstlisting}[language=lia]
!?[Titel](https://www.youtube
.com/watch?v=VIDEO_ID)

!?[Video](datei.mp4)
\end{lstlisting}
\tcblower
\prevsm{\faYoutube~YouTube, Vimeo und Videodateien werden direkt eingebettet.\\[2pt]
\textbf{!?} kombiniert Bild und Ton zu einem abspielbaren Video.}
\end{mdcard}

% --- R6. Eingebettete Ressource -------------------------------------
\begin{mdcard}{Eingebettete Ressource}
\begin{lstlisting}[language=lia]
??[Katalog](https://kxp.k10plus.de)

??[Spezifikation](
  https://dublincore.org/...)
\end{lstlisting}
\tcblower
\prevsm{\faWindowMaximize~PDF, GeoGebra, 3D-Modelle und Webseiten direkt anzeigen.\\[2pt]
\textbf{??} bettet die Ressource ein, statt sie nur zu verlinken.}
\end{mdcard}

% --- A8. Zitate & Hinweise -----------------------------------------
\begin{mdcard}{Zitate \& Hinweise}
\begin{lstlisting}[language=lia]
> Ein einfaches Zitat.

> [!NOTE]    blau
> [!TIP]     grün
> [!IMPORTANT] lila
> [!WARNING] orange
> [!CAUTION] rot
\end{lstlisting}
\tcblower
\prevtxt{%
\textcolor{rulegray}{\rule{1pt}{1.1em}}\hspace{4pt}\textit{Ein einfaches Zitat.}\\[2pt]
\textcolor{mdblue}{\rule{1pt}{1.0em}}\hspace{3pt}{\scriptsize\textbf{NOTE}\,/\,\textbf{TIP}\,/\,\textbf{IMPORTANT}}\\
\hspace*{8pt}{\scriptsize\textbf{WARNING}\,/\,\textbf{CAUTION}}\\
\hspace*{8pt}{\scriptsize Folgezeilen = Inhalt der Box.}}
\end{mdcard}

% =====================================================================
% Block D — Referenz-Karten (in den Aufgaben nicht zwingend gebraucht)
% =====================================================================

\end{multicols}
\endgroup

\sheetfooter

\newpage

% =====================================================================
% SEITE 2 — LIASCRIPT
% =====================================================================
\noindent\begin{tcolorbox}[
  enhanced, frame hidden, sharp corners,
  colback=liagreen, coltext=white,
  left=12pt, right=8pt, top=4pt, bottom=4pt,
  before skip=0pt, after skip=4pt,
  width=\linewidth,
  borderline west={4pt}{-4pt}{liagreen!55}
]
\begin{minipage}[c]{0.66\linewidth}
  {\sffamily\fontsize{26}{28}\selectfont\bfseries LiaScript}\hspace{0.4em}%
  {\sffamily\large\color{white!75} — die Erweiterungen}\\[2pt]
  {\sffamily\footnotesize\itshape\color{white!90} Alles, was Markdown allein nicht kann: Quiz, Animationen, ausführbarer Code, Sprache.}
\end{minipage}\hfill
\begin{minipage}[c]{0.20\linewidth}\raggedright
  {\sffamily\scriptsize\color{white!85}\faGraduationCap~~Fortbildung~$\cdot$~Freiberg 2026}\\[2pt]
  {\sffamily\bfseries\normalsize Seite 2\,/\,5}\\[1pt]
  {\sffamily\scriptsize\color{white!75} Erweiterungen \& Interaktion}
\end{minipage}\hfill
\begin{minipage}[c]{0.10\linewidth}\centering
  {\color{black}\colorbox{white}{\hspace{1pt}\qrcode[height=1.55cm,nolinks]{https://github.com/MINT-the-GAP}\hspace{1pt}}}\\[1pt]
  {\sffamily\scriptsize\color{white!90}\faGithub~MINT-the-GAP}
\end{minipage}
\end{tcolorbox}

% --- Tag für Aufgabenbezug (grün, oben rechts im Karten-Titel) -----
\newcommand{\aufglia}[1]{\hfill{\sffamily\scriptsize\color{white!85}\faTasks~A#1}}

\begin{multicols}{3}

% =====================================================================
% Block B — LiaScript-Erweiterungen (Aufgaben 9-11)
% =====================================================================

% --- A10. Animation / Click-Reveal ----------------------------------
\begin{liacard}{Animation \& Click-Reveal}
\begin{lstlisting}[language=lia]
Sofort sichtbar.

{{1}}
Erst nach einem Klick.

{{2}}
Nach dem zweiten Klick.
\end{lstlisting}
\tcblower
\prevsm{Sofort sichtbar.\\
{\color{rulegray}---}\\
\faMousePointer~nach 1. Klick: \textit{\dots}\\
\faMousePointer~nach 2. Klick: \textit{\dots}}
\end{liacard}

% --- A11. Sprecher-Text ---------------------------------------------
\begin{liacard}{Sprecher-Text (TTS)}
\begin{lstlisting}[language=lia]
## Folie
         --{{0}}--
Dieser Text wird vorgelesen,
sobald die Folie erscheint.
\end{lstlisting}
\tcblower
\prevsm{\faVolumeUp~Wird automatisch
gesprochen.\\
\texttt{\{\{0\}\}}~=~beim Öffnen,
\texttt{\{\{1\}\}}~=~nach 1.~Klick~\dots\\[2pt]
{\footnotesize Stimme im Header: \texttt{narrator:}}}
\end{liacard}

% =====================================================================
% Block C — Quizze (Aufgaben 12-13)
% =====================================================================

% --- A12. Freitext --------------------------------------------------
\begin{liacard}{Quiz: Freitext}
\begin{lstlisting}[language=lia]
Welcher Operator verkleinert
die Treffermenge?
[[AND]]
[[?]] Beide Begriffe müssen
      vorkommen.
\end{lstlisting}
\tcblower
\prevsm{Welcher Operator verkleinert die Treffermenge?\\[2pt]
\fbox{\strut\hspace{6em}}\\[3pt]
\faQuestionCircle~Hinweis (\texttt{[[?]]}, klappbar)\\[2pt]
{\footnotesize Mehrere gültige Antworten: weitere \texttt{[[\dots]]} anhängen.}}
\end{liacard}

% --- A13a. Single-Choice --------------------------------------------
\begin{liacard}{Quiz: Single-Choice}
\begin{lstlisting}[language=lia]
- [( )] falsch
- [(X)] richtig
- [[?]] Hinweis
*******************

Bravo, alles richtig!

*******************
\end{lstlisting}
\tcblower
\prevsm{$\circ$~falsch\\
$\bullet$~\textbf{richtig}\\
\faQuestionCircle~Hinweis (klappbar)\\[3pt]
{\footnotesize\color{liagreen}\faCheck~\textit{Bravo, alles richtig!}}\\[2pt]
{\footnotesize Lösungsblock nach \texttt{*****}\,---\,erscheint erst bei korrekter Antwort.}}
\end{liacard}

% --- A13b. Multiple-Choice ------------------------------------------
\begin{liacard}{Quiz: Multiple-Choice}
\begin{lstlisting}[language=lia]
- [[X]] vergrößert Treffer
- [[X]] verknüpft Synonyme
- [[ ]] schließt Begriffe aus
- [[X]] ist Vereinigungsmenge
\end{lstlisting}
\tcblower
\prevsm{\faCheckSquare~vergrößert Treffer\\
\faCheckSquare~verknüpft Synonyme\\
\faSquare~schließt Begriffe aus\\
\faCheckSquare~ist Vereinigungsmenge\\[2pt]
{\footnotesize \texttt{[[X]]} = beliebig viele richtig.}}
\end{liacard}

% --- A13c. Kachelquiz -----------------------------------------------
\begin{liacard}{Quiz: Kacheln}
\begin{lstlisting}[language=lia]
<!-- data-randomize="true" -->
Himmel: [->[rot|(blau)|grün]]

Gras: [->[(grün)|blau|rot]]
\end{lstlisting}
\tcblower
\prevsm{\faTh~Antwortkacheln in die passenden Lücken ziehen.\\[2pt]
Die richtige Lösung steht in \textbf{runden Klammern}; Alternativen trennt \texttt{|}.\\[2pt]
{\footnotesize LiaScript-Standard, kein Import nötig.}}
\end{liacard}

% =====================================================================
% Block D — Templates & Quiz-Kommentare
% =====================================================================

% --- JSXGraph --------------------------------------------------------
\begin{liacard}{JSXGraph}
\begin{lstlisting}[language=lia]
```javascript @JSX.Graph
board.create('point', [1, 2],
  {name: 'A'});
```
\end{lstlisting}
\tcblower
\prevsm{\faProjectDiagram~Interaktive Geometrie, Funktionen und Datenvisualisierung.\\[2pt]
Der JavaScript-Code erhält automatisch das Board \texttt{board}.}
\end{liacard}

% --- ABCjs -----------------------------------------------------------
\begin{liacard}{ABCjs: Musiknotation}
\begin{lstlisting}[language=lia]
```abc @ABCJS.render
X: 1
M: 4/4
L: 1/4
K: C
C D E F | G4 |
```
\end{lstlisting}
\tcblower
\prevsm{\faMusic~ABC-Notation als Noten darstellen und direkt abspielen.\\[2pt]
\texttt{@ABCJS.eval} macht den Notentext zusätzlich editierbar.}
\end{liacard}

% --- Algebrite -------------------------------------------------------
\begin{liacard}{Algebrite: Terme prüfen}
\begin{lstlisting}[language=lia]
Forme äquivalent um:
[[(x - 2) * (x + 2)]]
@Algebrite.check(`x^2 - 4`)

Ableitung: [[-(K/2) * t^(-3/2)]]
@Algebrite.check(`-(K/2) * t^(-3/2)`)

Ergebnis: [[1/3]]
@Algebrite.check2(1/3,0.01)

Schätzung: [[1200]]
@Algebrite.check_margin(1100,1400)
\end{lstlisting}
\tcblower
\prevsm{\faCalculator~Komplexe Terme in Backticks schützen; äquivalente Schreibweisen werden akzeptiert.\\[2pt]
Auch TeX-/LaTeX-Eingaben wie \texttt{\textbackslash frac\{1\}\{3\}} werden erkannt und vor der Prüfung umgewandelt.\\[2pt]
\texttt{check2} ergänzt eine Toleranz; \texttt{check\_margin} prüft einen Wertebereich.}
\end{liacard}

% --- Speech Recognition ---------------------------------------------
\begin{liacard}{Speech-Recognition-Quiz}
\begin{lstlisting}[language=lia]
<!-- data-solution-button="off" -->
Sprich: Hallo.
[[!]]
@SpeechRecognition(de-DE,Hallo)
\end{lstlisting}
\tcblower
\prevsm{\faMicrophone~Prüft eine gesprochene Phrase.\\[2pt]
\texttt{.withFeedback(...)} zeigt bei Abweichung das erkannte Transkript.\\[2pt]
{\footnotesize Benötigt Chrome, Edge oder Opera.}}
\end{liacard}

% --- Häufige Quiz-Kommentare ----------------------------------------
\begin{liacard}{Häufige Quiz-Kommentare}
\begin{lstlisting}[language=lia]
<!-- data-solution-button="5"
     data-hint-button="3"
     data-randomize="true"
     data-show-partial-solution=
  "true" -->
Frage ...
[[Antwort]]
\end{lstlisting}
\tcblower
\prevsm{\faSlidersH~In Aufgabensammlung und Wochenaufgaben regelmäßig genutzt.\\[2pt]
Lösung nach 5, Hinweis nach 3 Versuchen; mischen, Teilergebnisse anzeigen und Lösung zeitverzögert freigeben.}
\end{liacard}

\end{multicols}

% --- Untere Zeile: ausgewählte LiaScript-Templates (volle Breite) ---
\noindent
\begin{liainfo}{\faPuzzlePiece~~Vier LiaScript-Templates — jede benötigte Zeile in den Dokumentkopf schreiben}
\scriptsize
\begin{minipage}[t]{0.24\linewidth}
\textbf{\color{liagreen}\faProjectDiagram~~JSXGraph}\\[1pt]
\texttt{import:}\\[-1pt]
{\tiny\url{https://raw.githubusercontent.com/liaTemplates/JSXGraph/main/README.md}}\\[1pt]
Codeblock mit \texttt{@JSX.Graph}.
\end{minipage}\hfill
\begin{minipage}[t]{0.24\linewidth}
\textbf{\color{liagreen}\faMusic~~ABCjs}\\[1pt]
\texttt{import:}\\[-1pt]
{\tiny\url{https://raw.githubusercontent.com/liaTemplates/ABCjs/main/README.md}}\\[1pt]
ABC-Block mit \texttt{@ABCJS.render}.
\end{minipage}\hfill
\begin{minipage}[t]{0.24\linewidth}
\textbf{\color{liagreen}\faCalculator~~Algebrite}\\[1pt]
\texttt{import:}\\[-1pt]
{\tiny\url{https://raw.githubusercontent.com/liaTemplates/algebrite/master/README.md}}\\[1pt]
Nach Eingabe: \texttt{@Algebrite.check(...)}.
\end{minipage}\hfill
\begin{minipage}[t]{0.24\linewidth}
\textbf{\color{liagreen}\faMicrophone~~Speech Recognition}\\[1pt]
\texttt{import:}\\[-1pt]
{\tiny\url{https://raw.githubusercontent.com/LiaTemplates/Speech-Recognition-Quiz/refs/heads/main/README.md}}\\[1pt]
Generisches \texttt{[[!]]}-Quiz + Makro.
\end{minipage}
\end{liainfo}

\sheetfooter

\newpage

% =====================================================================
% SEITE 3 — SCHULLIA: UNTERRICHT & INTERAKTION
% =====================================================================
\noindent\begin{tcolorbox}[
  enhanced, frame hidden, sharp corners,
  colback=schulred, coltext=white,
  left=12pt, right=8pt, top=4pt, bottom=4pt,
  before skip=0pt, after skip=4pt,
  width=\linewidth,
  borderline west={4pt}{-4pt}{schulred!55}
]
\begin{minipage}[c]{0.66\linewidth}
  {\sffamily\fontsize{26}{28}\selectfont\bfseries SchulLia}\hspace{0.4em}%
  {\sffamily\large\color{white!75} — Unterricht \& Interaktion}\\[2pt]
  {\sffamily\footnotesize\itshape\color{white!90} Makros für Abgabe, Bewertung, Sprachaufgaben und lernförderliche Rückmeldung.}
\end{minipage}\hfill
\begin{minipage}[c]{0.20\linewidth}\raggedright
  {\sffamily\scriptsize\color{white!85}\faGraduationCap~~Fortbildung~$\cdot$~Freiberg 2026}\\[2pt]
  {\sffamily\bfseries\normalsize Seite 3\,/\,5}\\[1pt]
  {\sffamily\scriptsize\color{white!75} stabile Makros + Hinweise}
\end{minipage}\hfill
\begin{minipage}[c]{0.10\linewidth}\centering
  {\color{black}\colorbox{white}{\hspace{1pt}\qrcode[height=1.55cm,nolinks]{https://mint-the-gap.github.io/Aufgabensammlung/}\hspace{1pt}}}\\[1pt]
  {\sffamily\scriptsize\color{white!90}\faTasks~Aufgabensammlung}
\end{minipage}
\end{tcolorbox}

\begin{multicols}{3}

\begin{schulcard}{Abgabe \& Auswertung}{lia-freeze-v2/main/README.md}
\begin{lstlisting}[language=lia]
@Abgabe

@Auswertung(F12;Tab;Time)
\end{lstlisting}
\tcblower
\prevsm{\faShareSquare~Erzeugt den Abgabelink.\\[2pt]
\faChartBar~Auswertung mit Formular-ID, Tabellenblatt und Zeitspalte.}
\end{schulcard}

\begin{schulcard}{Prüfung \& Bewertung}{lia-freeze-v2/main/README.md}
\begin{lstlisting}[language=lia]
@Exam(60min)

3 + 4 = [[7]]

@ADetails(2;Addition,
          Grundrechenarten)
\end{lstlisting}
\tcblower
\prevsm{\faClock~Prüfungsdauer: 60 Minuten.\\[2pt]
\faTasks~\texttt{@ADetails} ordnet dem unmittelbar vorhergehenden Quiz 2 BE zu.\\[2pt]
\faTags~Nach dem Semikolon folgen optionale, kommagetrennte Themen-Tags. \texttt{@Auswertung} berechnet damit Gesamtpunktzahl und Themenbereiche.\\[2pt]
{\footnotesize Ohne Tags: \texttt{@ADetails(2)}; Standardwert: 1 BE. Leerzeile vor dem Makro.}}
\end{schulcard}

\begin{schulcard}{Text markieren}{lia-marker/main/README.md}
\begin{lstlisting}[language=lia]
<div class="markerquiz">
@markred(Nomen)
@markblue(Verb)
@TextmarkerQuiz
</div>

@markedgreen(Beispiel)
\end{lstlisting}
\tcblower
\prevsm{\textcolor{red}{\faHighlighter~Nomen}\quad
\textcolor{blue}{\faHighlighter~Verb}\\[3pt]
Zielwörter als Quiz in der erwarteten Farbe markieren; im Erklärmodus erscheint die Worterläuterung in der ausgewählten Sprache.\\[2pt]
{\footnotesize \texttt{@mark(...)} akzeptiert jede Farbe; \texttt{@markedCOLOR(...)} zeigt eine fertige Markierung.}}
\end{schulcard}

\begin{schulcard}{Kachelfolgen}{lia-kachel/main/README.md}
\begin{lstlisting}[language=lia]
@Kachelfolge(`[->[(A)|B]]
              [->[(C)|D]]`)

@KachelfolgeN(`[->[(1)|2]]
               [->[(2)|3]]`)
\end{lstlisting}
\tcblower
\prevsm{\faTh~Kacheln zu passenden Zielen ziehen.\\[2pt]
\texttt{@Kachelfolge}: Reihenfolge egal.\\
\texttt{@KachelfolgeN}: Ziele schrittweise aufdecken.\\[2pt]
{\footnotesize Backticks um die gesamte Folge.}}
\end{schulcard}

\begin{schulcard}{Flexible Spalten (DynFlex)}{lia-DynFlex/main/README.md}
\begin{lstlisting}[language=lia]
<section class="dynFlex">
<div class="flex-child">Inhalt A</div>
<div class="flex-child">Inhalt B</div>
</section>
\end{lstlisting}
\tcblower
\prevsm{\faColumns~Erzeugt nebeneinanderliegende, per Ziehen skalierbare Bereiche.\\[2pt]
Jeder Bereich braucht \texttt{class="flex-child"}; auf kleinen Displays stapeln sie sich.}
\end{schulcard}

\begin{schulcard}{Inhalte nach Ansicht}{lia-board-mode/main/README.md}
\begin{lstlisting}[language=lia]
<div data-lia-only="slides">
Nur auf Folien
</div>

<div data-lia-only="presentation">
Nur in der Präsentation
</div>

<div data-lia-only="textbook">
Nur im Lehrbuch
</div>
\end{lstlisting}
\tcblower
\prevsm{\faEye~\texttt{data-lia-only} zeigt einen Block ausschließlich in der angegebenen Ansicht.\\[2pt]
Unterstützt werden \texttt{slides}, \texttt{presentation} und \texttt{textbook}.\\[2pt]
{\footnotesize Das Template erweitert außerdem Folien und Präsentation auf nahezu volle Breite.}}
\end{schulcard}

\begin{schulcard}{Rechtschreibung \& Diktat}{lia-orthography/main/README.md}
\begin{lstlisting}[language=lia]
@orthography(
 `<!-- data-solution-button="2" -->`,
 `Feler im Saz.`,
 `Fehler im Satz.`)

@diktat(Schmetterling)
\end{lstlisting}
\tcblower
\prevsm{\faSpellCheck~Fehler im Ausgangstext finden und verbessern.\\[2pt]
\faVolumeUp~\texttt{@diktat} spricht Wort oder Satz vor; die Eingabe wird geprüft.\\[2pt]
{\footnotesize Reihenfolge: Quiz-Kommentar, fehlerhafter Text, korrekte Lösung.}}
\end{schulcard}

\begin{schulcard}{Bruchteile markieren}{lia-Mathe/main/README.md}
\begin{lstlisting}[language=lia]
@circleQuiz(3/8)

@rectQuiz(1/3)

@circleQuizC(3/8,
 `<!-- data-hint-button="2" -->`)
\end{lstlisting}
\tcblower
\prevsm{\faCircleNotch~Drei von acht Kreissektoren auswählen.\\[2pt]
\faThLarge~Ein Drittel eines Rechtecks markieren.\\[2pt]
{\footnotesize \texttt{...C} übernimmt Quiz-Kommentare.}}
\end{schulcard}

\begin{schulcard}{Zeitgesteuerte Lösung}{lia-timer/main/README.md}
\begin{lstlisting}[language=lia]
<!-- data-solution-timer="15s"
     data-solution-timer-start=
     "oncheck"
     data-solution-timer-badge=
     "off" -->
7 + 8 = [[15]]
\end{lstlisting}
\tcblower
\prevsm{\faHourglassHalf~Lösung erst nach 15 Sekunden freigeben.\\[2pt]
Start: \texttt{immediate}, \texttt{oncheck} oder \texttt{onclick}.}
\end{schulcard}

\begin{schulcard}{Globale Quizregeln}{lia-globalquiz/main/README.md}
\begin{lstlisting}[language=lia]
@global(data-solution-button=3
        data-hint-button=2)

@globalforced(
  data-hint-button=4)
\end{lstlisting}
\tcblower
\prevsm{\faSlidersH~Vorgaben für alle folgenden Quizze setzen.\\[2pt]
\texttt{@global}: lokale Werte dürfen abweichen.\\
\texttt{@globalforced}: lokale Werte überschreiben.}
\end{schulcard}

\begin{schulcard}{Quiz einzeln zurücksetzen}{lia-resetter/main/README.md}
\begin{lstlisting}[language=lia]
7 + 8 = [[15]]

@resetter
\end{lstlisting}
\tcblower
\prevsm{\faUndo~Direkt nach dem Quiz einfügen.\\[2pt]
Setzt nur dieses Quiz einschließlich Eingabe und Prüfzustand vollständig zurück.}
\end{schulcard}

\begin{schulcard}{Handschrift \& OCR}{lia-canvas-ocr/main/README.md}
\begin{lstlisting}[language=lia]
Berechne:
$3x + 2 = 11$
[[x = 3]]
@canvas
\end{lstlisting}
\tcblower
\prevsm{\faPen~\texttt{@canvas} öffnet direkt am Antwortfeld eine Schreibfläche.\\[2pt]
Der markierte handschriftliche Bereich wird erkannt und als LaTeX in das Eingabefeld übertragen.}
\end{schulcard}

\begin{schulcard}{Annotationen}{lia-annotation/main/README.md}
\begin{lstlisting}[language=lia]
# Tafelbild

Skizze und Rechnung
\end{lstlisting}
\tcblower
\prevsm{\faPen~Zeichenleiste auch über JSXGraph-/\texttt{lia-coordinate}-Inhalten: DGS-Konstruktionen bearbeiten und Koordinatensysteme skizzieren.\\[2pt]
{\footnotesize Benötigt beide Imports; OCR ist enthalten.}}
\end{schulcard}

\end{multicols}

\sheetfooter

\newpage

% =====================================================================
% SEITE 4 — SCHULLIA: FÄCHER & VISUALISIERUNG
% =====================================================================
\noindent\begin{tcolorbox}[
  enhanced, frame hidden, sharp corners,
  colback=schulred, coltext=white,
  left=12pt, right=8pt, top=4pt, bottom=4pt,
  before skip=0pt, after skip=4pt,
  width=\linewidth,
  borderline west={4pt}{-4pt}{schulred!55}
]
\begin{minipage}[c]{0.66\linewidth}
  {\sffamily\fontsize{26}{28}\selectfont\bfseries SchulLia}\hspace{0.4em}%
  {\sffamily\large\color{white!75} — Fächer \& Visualisierung}\\[2pt]
  {\sffamily\footnotesize\itshape\color{white!90} Freie Antworten mit KI, Koordinatensysteme, Funktionen und dynamische Geometrie.}
\end{minipage}\hfill
\begin{minipage}[c]{0.20\linewidth}\raggedright
  {\sffamily\scriptsize\color{white!85}\faGraduationCap~~Fortbildung~$\cdot$~Freiberg 2026}\\[2pt]
  {\sffamily\bfseries\normalsize Seite 4\,/\,5}\\[1pt]
  {\sffamily\scriptsize\color{white!75} Mathematik + Templates}
\end{minipage}\hfill
\begin{minipage}[c]{0.10\linewidth}\centering
  {\color{black}\colorbox{white}{\hspace{1pt}\qrcode[height=1.55cm,nolinks]{https://wochenaufgaben.gsg-freiberg.de/}\hspace{1pt}}}\\[1pt]
  {\sffamily\scriptsize\color{white!90}\faTasks~Wochenaufgaben}
\end{minipage}
\end{tcolorbox}

\begingroup
\setlength{\multicolsep}{2pt}
\tcbset{cheatbase/.append style={top=1pt,bottom=1pt,before skip=2pt,after skip=2pt,fontlower=\footnotesize}}
\begin{multicols}{3}

\begin{schulcard}{LLM-Quiz: freie Antworten}{lia-llm/main/README.md}
{\sffamily\scriptsize\bfseries A: Vollständige Lösungsalternativen}\par\vspace{2pt}
\begin{lstlisting}[language=lia,breakatwhitespace=false,breakautoindent=false]
Erkläre die Fotosynthese.

<!-- data-solution-button="off"
     data-llm-textarea="3" -->
[[Antwort]]
```text @LLMQuiz(0.66;assessmentengine=quality;operator=erklaeren;solution=1;feedback=1;Rechtschreibung=1;Satzbau=1;maxthinkingtime=15s;maxthinkingtokens=medium,`Erkläre die Fotosynthese.`)
Pflanzen bilden mit Licht aus
Wasser und Kohlenstoffdioxid
Glucose. Dabei entsteht Sauerstoff.
<!-- lia-llm:alternative -->
Mit Lichtenergie bauen Pflanzen
aus Wasser und CO2 Traubenzucker
auf und geben Sauerstoff ab.
```
\end{lstlisting}

\par\vspace{5pt}
{\sffamily\scriptsize\bfseries B: Einzelne Kriterien}\par\vspace{2pt}
\begin{lstlisting}[language=lia,breakatwhitespace=false,breakautoindent=false]
Nenne beide Ausgangsstoffe
der Fotosynthese.

<!-- data-solution-button="off" -->
[[Antwort]]
```text @LLMQuiz(0.55;coverage=1;solution=1;feedback=1,`Nenne beide Ausgangsstoffe der Fotosynthese.`)
<!-- lia-llm:criterion -->
Wasser ist ein Ausgangsstoff.
<!-- lia-llm:criterion -->
CO2 ist ein Ausgangsstoff.
<!-- lia-llm:solution -->
Die Ausgangsstoffe sind Wasser
und Kohlenstoffdioxid.
```
\end{lstlisting}
\tcblower
\prevsm{\faComments~\textbf{A:} \texttt{alternative} trennt vollständige Musterlösungen: Eine davon muss sinngemäß erfüllt sein.\\[3pt]
\texttt{quality} prüft gründlicher; \texttt{operator=erklaeren} verlangt eine Erklärung.\\[3pt]
\texttt{Rechtschreibung=1} und \texttt{Satzbau=1}: nach der Inhaltsprüfung auf \textbf{Sprache prüfen} klicken.\\[3pt]
\texttt{15s} und \texttt{medium} (512 Tokens) begrenzen zusätzliches Denken.\\[3pt]
\faTasks~\textbf{B:} Jeder \texttt{criterion}-Marker beginnt eine Anforderung. \texttt{coverage=1} verlangt beide; \texttt{solution} leitet die Musterlösung ein.\\[3pt]
\texttt{0.66}/\texttt{0.55}: Prüfschwelle. \texttt{solution=1}: Lösung nach richtiger Antwort. \texttt{feedback=1}: Rückmeldung.\\[3pt]
\textbf{A und B getrennt verwenden:} Kriterien nicht mit \texttt{alternative} oder \texttt{operator} mischen.\\[3pt]
Jeweils den gesamten Makroaufruf in einer Quelltextzeile schreiben; zweiter Parameter ist die vollständige Frage.}
\end{schulcard}

\begin{schulcard}{Gleichungslücke}{lia-Mathe/Proposals/README.md}
\begin{lstlisting}[language=lia]
$\frac{\liaquiz}{3}
 = \frac{5}{15}$
@liaQuiz(1)

$x + \liaquiz = 7$
@liaQuizC(3,`<!-- ... -->`)
\end{lstlisting}
\tcblower
\prevsm{\faFlask~\texttt{\textbackslash liaquiz} setzt ein Eingabefeld direkt in eine Gleichung.\\[2pt]
Das folgende Makro liefert die Lösung.}
\end{schulcard}

\begin{schulcard}{Koordinatensystem}{lia-coordinate/main/README.md}
\begin{lstlisting}[language=lia]
@Koordinatensystem(`xmin=-5;
  xmax=5;ymin=-4;ymax=4;
  width=800;id=B`)

@AchsenBeschriftung(`id=B;
  xlabel=$x$;ylabel=$y$`)
\end{lstlisting}
\tcblower
\prevsm{\faBorderAll~System \texttt{B} mit Wertebereich und Breite anlegen.\\[2pt]
\faProjectDiagram~Achsen anschließend beschriften.\\[2pt]
{\footnotesize Zuerst JSXGraph, danach \texttt{lia-coordinate} separat importieren.}}
\end{schulcard}

\begin{schulcard}{DGS: bauen \& exportieren}{lia-coordinate/main/README.md}
\begin{lstlisting}[language=lia]
@DGS(`B;tools=[200;310;700];
  restrictions=[200;300]`)

@Zirkel(`B`)

@Geodreieck(`B`)
\end{lstlisting}
\tcblower
\prevsm{\faDraftingCompass~\textbf{1. Bauen:} Im DGS Punkte, Geraden, Funktionen und Konstruktionen anlegen.\\[2pt]
\faDownload~\textbf{2. Exportieren:} Objektliste öffnen und den erzeugten Makroblock kopieren.\\[2pt]
\faEdit~\textbf{3. Einfügen:} Makroblock im LiveEditor einsetzen und weiterbearbeiten.}
\end{schulcard}

\begin{schulcard}{Flächen, Winkel \& Kreise}{lia-coordinate/main/README.md}
\begin{lstlisting}[language=lia]
@Flaeche(`B;[A;P;C];#e9c46a;
  0.25;inhalt=1;umfang=1`)

@Winkel(`B;alpha;[A;P;C];
  #e76f51;0.8;Wert=1`)

@Kreis(`B;k;M;#264653;0.2;
  radius=P;inhalt=1`)

@Kreissektor(`B;[M;P;Q];#e9c46a;
  0.25;s;inhalt=1;umfang=1`)
\end{lstlisting}
\tcblower
\prevsm{\faShapes~Polygonfläche, Winkel, Kreis und Kreisbogen aus benannten Punkten erzeugen.\\[2pt]
\texttt{@Kreissektor} verwendet Mittelpunkt, Radiuspunkt und zweiten Schenkelpunkt; \texttt{@Kreissegment} ist ein Alias.}
\end{schulcard}

\begin{schulcard}{Analysepunkte}{lia-coordinate/main/README.md}
\begin{lstlisting}[language=lia]
@Nullstellen(`B;f;blue;N;wert=1`)

@Extrempunkte(`B;f;green;E;wert=1`)

@Wendepunkte(`B;f;magenta;W;wert=1`)

@Schnittpunkt(`B;f;g;orange;S;wert=1`)
\end{lstlisting}
\tcblower
\prevsm{\faMapMarker~Nullstellen, Extrem- und Wendepunkte einer benannten Funktion dynamisch bestimmen.\\[2pt]
\texttt{@Schnittpunkt} markiert Schnittpunkte zweier Funktionen, Geraden oder Kreise.}
\end{schulcard}

\begin{schulcard}{Funktionen \& Werkzeuge}{lia-coordinate/main/README.md}
\begin{lstlisting}[language=lia]
@PlotFunktion(`B;f;0.5*x^2-2;red`)

@Tangente(`B;f;[2;0];magenta;t;T`)

@Tabelle(`n=3;x;f;P;id=B`)

@Regression(`B`)

@Rekonstruktion(`B;2x-1;0.1`)
\end{lstlisting}
\tcblower
\prevsm{\faChartLine~Funktion zeichnen, Tangente im Berührpunkt anlegen und eine Wertetabelle ergänzen.\\[2pt]
\texttt{@Regression} öffnet Zeichen- und Ausgleichswerkzeuge; \texttt{@Rekonstruktion} prüft den Zielterm mit Toleranz.}
\end{schulcard}

\begin{schulcard}{Punkte \& Geraden}{lia-coordinate/main/README.md}
\begin{lstlisting}[language=lia]
@Punkt(`B;A;2;3;#e63946;1;fix`)

@ErzeugePunkt(`B;P;1;2`,`<!-- -->`)

@Gerade(`B;[A;P];#457b9d;g`)

@Vektor(`B;[A;P];#2a9d8f;v`)
\end{lstlisting}
\tcblower
\prevsm{\faMapMarker~Punkte anzeigen oder interaktiv erzeugen lassen.\\[2pt]
\faProjectDiagram~Geraden und Vektoren referenzieren vorhandene Punkte über ihre Namen.}
\end{schulcard}

\begin{schulcard}{Geometriequizze}{lia-coordinate/main/README.md}
\begin{lstlisting}[language=lia]
@UmfangQuiz(`B;3;12;0.05`,
 `<!-- data-solution-button="5" -->`)

@FlaecheQuiz(`B;3;6;0.05`,
 `<!-- -->`)

@KonstruktionQuiz(
 `B;3;fest;S4,W90,S3`,
 `<!-- data-solution-button="5" -->`)
\end{lstlisting}
\tcblower
\prevsm{\faRulerCombined~Prüft ein selbst konstruiertes Polygon mit Eckenzahl, Zielwert und absoluter Toleranz.\\[2pt]
\texttt{@UmfangQuiz}: Umfang 12; \texttt{@FlaecheQuiz}: Fläche 6.\\[2pt]
\texttt{@KonstruktionQuiz} prüft Seiten und Winkel in fester oder offener Anordnung.}
\end{schulcard}

\begin{schulcard}{Erklärhinweis \& Tooltip}{lia-mathpath/master/README.md}
\begin{lstlisting}[language=lia]
[[?]] @Explain

<div class="notip">
Ohne Glossar-Tooltip
</div>
\end{lstlisting}
\tcblower
\prevsm{\faLightbulb~\texttt{@Explain} erzeugt einen gestuften Lösungshinweis.\\[2pt]
Die Themen stammen aus \texttt{@ADetails(...)} beim Quiz; passende Erklärlinks werden in \texttt{Explain.md} zugeordnet.\\[2pt]
\texttt{notip} unterdrückt automatische Glossar-Tooltips im umschlossenen Bereich.}
\end{schulcard}

\end{multicols}
\endgroup

\noindent
\begin{schulinfo}{\faDownload~~Template-Imports}
\footnotesize
\begin{minipage}[t]{0.51\linewidth}
\textbf{Basis-URL:}  \texttt{https://raw.githubusercontent.com/MINT-the-GAP/} \\[1em]
Den Pfad oben rechts anfügen und die vollständige URL im LiaScript-Kopf unter \texttt{import:} eintragen.
\end{minipage}\hfill
\begin{minipage}[t]{0.48\linewidth}
\faInfoCircle~\textbf{Coordinate:} zuerst \texttt{https://raw.githubusercontent.com/liaTemplates/JSXGraph/main/README.md} importieren.\\[1pt]
\faClock~\texttt{lia-resetter/main/README.md}: Veröffentlichung folgt. 
\end{minipage}
\end{schulinfo}

\sheetfooter

\newpage

% =====================================================================
% SEITE 5 — GAMIFICATION MIT LIA-LOOT
% =====================================================================
\noindent\begin{tcolorbox}[
  enhanced, frame hidden, sharp corners,
  colback=lootteal, coltext=white,
  left=12pt, right=8pt, top=4pt, bottom=4pt,
  before skip=0pt, after skip=4pt,
  width=\linewidth,
  borderline west={4pt}{-4pt}{lootteal!55}
]
\begin{minipage}[c]{0.66\linewidth}
  {\sffamily\fontsize{26}{28}\selectfont\bfseries Gamification}\hspace{0.4em}%
  {\sffamily\large\color{white!75} — mit lia-loot}\\[2pt]
  {\sffamily\footnotesize\itshape\color{white!90} Ressourcen, Funde, Werkzeuge, Tore und Geheimnisse — die vollständige öffentliche Makrooberfläche.}
\end{minipage}\hfill
\begin{minipage}[c]{0.20\linewidth}\raggedright
  {\sffamily\scriptsize\color{white!85}\faGraduationCap~~Fortbildung~$\cdot$~Freiberg 2026}\\[2pt]
  {\sffamily\bfseries\normalsize Seite 5\,/\,5}\\[1pt]
  {\sffamily\scriptsize\color{white!75} alle öffentlichen Makros}
\end{minipage}\hfill
\begin{minipage}[c]{0.10\linewidth}\centering
  {\color{black}\colorbox{white}{\hspace{1pt}\qrcode[height=1.55cm,nolinks]{https://github.com/MINT-the-GAP/lia-loot}\hspace{1pt}}}\\[1pt]
  {\sffamily\scriptsize\color{white!90}\faGithub~lia-loot}
\end{minipage}
\end{tcolorbox}

\begingroup
\setlength{\multicolsep}{2pt}
\tcbset{cheatbase/.append style={top=1pt,bottom=1pt,before skip=2pt,after skip=2pt}}
\begin{multicols}{3}

\begin{lootcard}{Start: Ressourcen, Highscore, Erfolge}
\begin{lstlisting}[language=lia]
@Ressourcen(10, 3, 5)

@Highscore(100, 10, 5,
           15, 2)

@achievements
\end{lstlisting}
\tcblower
\prevsm{\faSlidersH~\textbf{Ressourcen:} \lootval{10} Gold · \lootval{3} Diamanten · \lootval{5} Energie.\\[2pt]
\faChartLine~\textbf{Highscore:} \lootval{100} Maximum · \lootval{10} Abzug je Fehlprüfung · \lootval{5} je geöffnetem Hinweis · \lootval{15} Freiminuten · danach \lootval{2}/Min. (sekundengenau).\\[2pt]
\faCheck~Aktiviert zwölf feste Erfolge. Aliasse: \texttt{@Achievements}, \texttt{@Erfolge}.}
\end{lootcard}

\begin{lootcard}{Gold, Diamanten \,\& Energie}
\begin{lstlisting}[language=lia]
@Schatztruhe(3; toc; menu)

@Diamanttruhe(2; anker)

@Energiekiste(3; mode)
\end{lstlisting}
\tcblower
\prevsm{\textbf{Gold:} \lootval{3} je in \texttt{toc} und \texttt{menu}.\\[1.5pt]
\textbf{Diamanten:} \lootval{2} inline; \texttt{anker} bindet den Fund an die Folie.\\[1.5pt]
\textbf{Energie:} \lootval{3} in \texttt{mode} (dritter Ressourcenwert nötig). Erster Wert = Menge; jedes Ziel erhält sie vollständig.}
\end{lootcard}

\begin{lootcard}{Schlüssel \,\& Schlösser}
\begin{lstlisting}[language=lia]
@Schluessel(blau; menu)

@Schloss(check, blau)

@Schloss(toc, rot; anker)
\end{lstlisting}
\tcblower
\prevsm{\faKey~\texttt{blau} = Schlüsselfarbe; \texttt{menu} = Fundort.\\[2pt]
\faPuzzlePiece~Beim Schloss steht erst das Ziel (\texttt{check}/\texttt{toc}), dann die nötige Farbe. Öffnen verbraucht einen passenden Schlüssel.\\[2pt]
{\footnotesize \texttt{anker} bindet ein sonst globales Schloss an seine Folie.}}
\end{lootcard}

\begin{lootcard}{Gemeinsame Fundoptionen}
\begin{lstlisting}[language=lia]
anker; 12s
theme=rot
farbmodus=dunkel
annotationen=aus
unsichtbar
erde; pflanze-zauberstaub
\end{lstlisting}
\tcblower
\prevsm{\texttt{anker} = Folienbindung; \lootval{12s} = zwölf Sekunden Wartezeit. Weitere Optionen steuern Theme, Farbmodus, Annotationen und Verbergung.\\[2pt]
Direkte Schichten \texttt{erde}/\texttt{pflanze} laufen von außen nach innen; Suffixe: \texttt{-unsichtbar}, \texttt{-zauberstaub}.}
\end{lootcard}

\columnbreak

\begin{lootcard}{Lupe \,\& Verbergung}
\begin{lstlisting}[language=lia]
@Lupe

@Unsichtbar(Geheimer Text)

@Zauberstaub(Hier funkelt etwas)
\end{lstlisting}
\tcblower
\prevsm{\faSearch~Lupe einsammeln und in der Leiste aktivieren.\\[2pt]
\texttt{@Unsichtbar} verdeckt vollständig, \texttt{@Zauberstaub} hinterlässt ein schwaches Funkeln.\\[2pt]
{\footnotesize Einzeiliger Inhalt; Spieleffekt, kein Zugriffsschutz.}}
\end{lootcard}

\begin{lootcard}{Aktionswerkzeuge}
\begin{lstlisting}[language=lia]
@Schaufel

@Giesskanne
\end{lstlisting}
\tcblower
\prevsm{\faPuzzlePiece~Schaufel öffnet Erde, Gießkanne bringt Pflanzen zum Blühen.\\[2pt]
Werkzeug immer vor der ersten Pflichtschicht erreichbar machen. Nur ein Aktionswerkzeug ist aktiv; die Lupe kann parallel laufen.}
\end{lootcard}

\begin{lootcard}{Erde: Block \,\& Inline}
\begin{lstlisting}[language=lia]
@Erdhaufen(unsichtbar)
Vergrabener Inhalt
@EndeErdhaufen

Text @Erdhaufen.inline(Hinweis)
\end{lstlisting}
\tcblower
\prevsm{\faEye~Blockform umschließt beliebigen LiaScript-Inhalt; Start und Ende stehen allein auf derselben Folie.\\[2pt]
\texttt{.inline} gibt einzeiligen Inhalt frei und benötigt kein Endmakro.}
\end{lootcard}

\begin{lootcard}{Pflanze: Block \,\& Inline}
\begin{lstlisting}[language=lia]
@Pflanze(zauberstaub)
Verborgener Inhalt
@EndePflanze

Text @Pflanze.inline(Hinweis)
\end{lstlisting}
\tcblower
\prevsm{Erst gießen, dann die Blüte anklicken: Danach wird der Inhalt sichtbar.\\[2pt]
Aliasse: \texttt{@Blume}, \texttt{@EndeBlume} und \texttt{@Blume.inline}.}
\end{lootcard}

\columnbreak

\begin{lootcard}{Puzzleteile \,\& Puzzletore}
\begin{lstlisting}[language=lia]
@Puzzleteil(rot; 1)
@Puzzleteil(rot; 2)
@Puzzletor(rot; [[2;1]])

@Puzzleteil(blau; 1)
@Puzzletor(blau; [[1]]; anker)
\end{lstlisting}
\tcblower
\prevsm{\faPuzzlePiece~\lootval{1}/\lootval{2} = Teilnummern; \texttt{[[2;1]]} = Teil 2 links, Teil 1 rechts. Alle Teile stehen davor; maximal 16.\\[2pt]
Ohne \texttt{anker}: Navigationstor. Mit \texttt{anker}: Inhalt bis zur nächsten \texttt{\#}-Überschrift.}
\end{lootcard}

\begin{lootcard}{Zweiweg- \,\& Einwegportale}
\begin{lstlisting}[language=lia]
@Portal(5)
@Portal(5; einweg)

@Einwegportal(5)
@Einbahnportal(5)
\end{lstlisting}
\tcblower
\prevsm{\faUndo~\lootval{5} = Ziel ist die fünfte Kursfolie (1-basiert). \texttt{@Portal} ist standardmäßig ein Zweiwegportal mit Rückweg.\\[2pt]
Die letzten drei Formen sind einweg; \texttt{@Einbahnportal} ist ein Alias.}
\end{lootcard}

\begin{lootcard}{Bedingte Bereiche}
\begin{lstlisting}[language=lia]
@lootif(Gold >= 5; spawn)
Dieser Inhalt erscheint dauerhaft.
@Endelootif
\end{lstlisting}
\tcblower
\prevsm{\faEye~\lootval{5} = Schwelle: Der Bereich erscheint ab fünf Gold dauerhaft. Einzig gültige Aktion: \texttt{spawn}.\\[2pt]
End-Aliasse: \texttt{@EndeLootif}, \texttt{@endlootif}, \texttt{@EndLootIf}.\\[2pt]
{\footnotesize Start und Ende allein, foliengleich und korrekt verschachtelt.}}
\end{lootcard}

\begin{lootcard}{Geheimfolie}
\begin{lstlisting}[language=lia]
## Das geheime Labor

@Geheimfolie
\end{lstlisting}
\tcblower
\prevsm{\faSearch~Verbirgt die Folie vor ToC und normaler Navigation. Erreichbar per exakter Titelsuche oder gezieltem Portal.\\[2pt]
{\footnotesize Kein Schutz für vertrauliche Inhalte: Quelle und Seitennummer bleiben prinzipiell einsehbar.}}
\end{lootcard}

\end{multicols}
\endgroup

\noindent
\begin{tcolorbox}[
  enhanced, colback=loottealLt, colframe=lootteal!75,
  boxrule=0.3pt, arc=2pt,
  left=5pt,right=5pt,top=1pt,bottom=1pt,
  before skip=1pt, after skip=0pt
]
\tiny
\begin{minipage}[t]{0.54\linewidth}
\textbf{\color{lootteal}\faDownload~~Import:}
\url{https://raw.githubusercontent.com/MINT-the-GAP/lia-loot/main/README.md}\\[-1pt]
Nur öffentliche Makros verwenden — niemals interne \texttt{@Loot...\_}-Aufrufe.
\end{minipage}\hfill
\begin{minipage}[t]{0.45\linewidth}
\textbf{Kosten:} Hinweis 1 Gold · Auflösen 1 Diamant · Prüfen 1 Energie.\quad
\textbf{Ziele:} \texttt{toc}, \texttt{menu}, \texttt{classroom}, \texttt{info}, \texttt{translator}, \texttt{mode}.\\[-1pt]
\textbf{Schloss zusätzlich:} \texttt{seitenwechsel}, \texttt{check}, \texttt{resolve}, \texttt{hint}, \texttt{portal}.\quad
\textbf{Trigger:} Aufgaben, Ressourcen/Kisten, Schloss/Puzzle, Geheimfolie/Lupe/Marker.
\end{minipage}
\end{tcolorbox}

\sheetfooter

\end{document}
